% ─────────────────────────────────────────────────────────────────────────────
%  Capítulo 6 – Backend x86-64
% ─────────────────────────────────────────────────────────────────────────────
\chapter{Backend x86-64}
\label{chap:backend}

\section{Convención de llamada System V AMD64 ABI}

El compilador respeta la convención de llamada estándar de Linux/macOS
para x86-64 (\emph{System V AMD64 ABI}):

\begin{table}[H]
  \centering
  \caption{Registros de la System V AMD64 ABI}
  \label{tab:abi}
  \begin{tabular}{lll}
    \toprule
    \textbf{Propósito}     & \textbf{Enteros}   & \textbf{Flotantes} \\
    \midrule
    Argumentos             & \reg{rdi}, \reg{rsi}, \reg{rdx}, \reg{rcx}, \reg{r8}, \reg{r9}
                           & \reg{xmm0}–\reg{xmm7} \\
    Valor de retorno       & \reg{rax}           & \reg{xmm0} \\
    Callee-saved           & \reg{rbx}, \reg{r12}–\reg{r15}, \reg{rbp} & \reg{xmm8}–\reg{xmm15} \\
    Caller-saved           & \reg{rcx}, \reg{rdx}, \reg{rsi}, \reg{rdi}, \reg{r8}–\reg{r11}
                           & \reg{xmm0}–\reg{xmm7} \\
    \bottomrule
  \end{tabular}
\end{table}

\section{Marco de pila}

Cada función establece su propio marco de pila siguiendo el
prólogo/epílogo estándar:

\begin{lstlisting}[style=asm, caption={Prólogo y epílogo de función}]
    pushq   %rbp
    movq    %rsp, %rbp
    subq    $N, %rsp        # N = frameSize (alineado a 16 bytes)
    ...
    movq    %rbp, %rsp
    popq    %rbp
    ret
\end{lstlisting}

Las variables locales y temporales se ubican en la pila según la fórmula:
\[
  \texttt{slot}(v) = -((\text{slot}_v + 1) \times 8)\texttt{(\%rbp)}
\]
donde $\text{slot}_v$ es el índice (0-based) asignado a la primera vez
que se usa la dirección de la variable en la IR.

\section{Generación de código por opcode}

\subsection{Operaciones aritméticas enteras}

\begin{table}[H]
  \centering
  \caption{Traducción de operaciones aritméticas enteras}
  \label{tab:arith}
  \begin{tabular}{lp{8cm}}
    \toprule
    \textbf{IR}     & \textbf{Secuencia x86-64 AT\&T} \\
    \midrule
    $t \leftarrow a + b$ &
      \texttt{movq mem(a), \%rax; addq mem(b), \%rax; movq \%rax, mem(t)} \\
    $t \leftarrow a * b$ &
      \texttt{movq mem(a), \%rax; imulq mem(b); movq \%rax, mem(t)} \\
    $t \leftarrow a / b$ &
      \texttt{movq mem(a), \%rax; cqto; idivq mem(b); movq \%rax, mem(t)} \\
    $t \leftarrow a \% b$ &
      \texttt{movq mem(a), \%rax; cqto; idivq mem(b); movq \%rdx, mem(t)} \\
    \bottomrule
  \end{tabular}
\end{table}

\subsection{Saltos condicionales}

Los saltos se implementan con el par \texttt{cmpq} + instrucción de salto
condicional. Para comparaciones que producen un booleano en la IR (p.ej.
\texttt{IR\_LT}), el patrón es:

\begin{lstlisting}[style=asm, caption={Comparación IR\_LT y salto condicional}]
    movq    mem(a), %rax
    cmpq    mem(b), %rax
    setl    %al              # 1 si a < b, 0 si no
    movzbq  %al, %rax
    movq    %rax, mem(t)
    ...
    testq   mem(cond), mem(cond)
    je      L_else           # JZ: salta si cond == 0
\end{lstlisting}

\subsection{Llamadas a función}

\begin{lstlisting}[style=asm, caption={Llamada a función y recogida de valor de retorno}]
    call    absValue         # llama a absValue()
    movq    %rax, mem(t)    # t = valor de retorno (%rax)
\end{lstlisting}

\subsection{E/S con \texttt{printf} y \texttt{scanf}}

\begin{lstlisting}[style=asm, caption={Impresión de entero con printf}]
    .section .rodata
fmt_int:    .string "%ld\n"
    .text
    leaq    fmt_int(%rip), %rdi
    movq    mem(val), %rsi
    xorl    %eax, %eax
    call    printf@PLT
\end{lstlisting}

\begin{lstlisting}[style=asm, caption={Lectura de entero con scanf}]
fmt_scanf_int: .string "%ld"
    leaq    fmt_scanf_int(%rip), %rdi
    leaq    mem(t)(%rbp), %rsi       # address of local variable in stack
    xorl    %eax, %eax
    call    scanf@PLT
\end{lstlisting}

\section{Ejemplo completo de programa compilado}

\begin{lstlisting}[style=java, caption={Programa fuente: suma hasta 5}]
public class TestWhile {
    public static int sumToFive() {
        int i; int s;
        i = 1; s = 0;
        while (i <= 5) { s = s + i; i = i + 1; }
        return s;
    }
    public static void main(String[] args) {
        int val;
        val = sumToFive();
        System.out.println(val);
    }
}
\end{lstlisting}

El ensamblador generado para \texttt{sumToFive} tiene la forma:

\begin{lstlisting}[style=asm, caption={Ensamblador x86-64 generado para sumToFive}]
    .globl  sumToFive
sumToFive:
    pushq   %rbp
    movq    %rsp, %rbp
    subq    $48, %rsp
L0:
    movq    -8(%rbp), %rax      # i
    cmpq    $5, %rax
    setle   %al
    movzbq  %al, %rax
    movq    %rax, -24(%rbp)     # t0 = i <= 5
    testq   -24(%rbp), %rax
    je      L1
    movq    -8(%rbp), %rax      # s = s + i
    addq    -16(%rbp), %rax
    movq    %rax, -32(%rbp)
    movq    -32(%rbp), %rax
    movq    %rax, -16(%rbp)
    movq    -8(%rbp), %rax      # i = i + 1
    addq    $1, %rax
    ...
    jmp     L0
L1:
    movq    -16(%rbp), %rax     # return s
    movq    %rbp, %rsp
    popq    %rbp
    ret
\end{lstlisting}
